\documentclass[conference]{IEEEtran}

\usepackage{iftex}
\ifPDFTeX
\usepackage[utf8]{inputenc}
\usepackage[T1]{fontenc}
\fi
\usepackage{amsthm}
\usepackage{amsmath,bm,amssymb}
\usepackage{newtxtext,newtxmath}
\interdisplaylinepenalty=2500
\usepackage{xcolor}
\usepackage{graphicx}
\usepackage{array,booktabs,adjustbox}
\usepackage{mathtools}
\usepackage{siunitx}
\usepackage[final]{microtype}
\usepackage{cite}
\usepackage[breaklinks=true,hidelinks]{hyperref}
\usepackage[nameinlink,capitalise,noabbrev]{cleveref}

\DeclareRobustCommand{\vct}[1]{\bm{#1}}
\DeclareRobustCommand{\mat}[1]{\bm{#1}}
\newcommand{\R}{\mathbb{R}}
\newcommand{\C}{\mathbb{C}}
\newcommand{\Hr}{^{\mathsf H}}
\newcommand{\T}{^{\mathsf T}}
\sisetup{detect-weight=true,detect-family=true,per-mode=symbol}
\emergencystretch=2em
\allowdisplaybreaks

\newtheoremstyle{bodysmall}{3pt}{3pt}{\normalfont\small}{0pt}{\bfseries\small}{.}{0.5em}{}
\theoremstyle{bodysmall}
\newtheorem{proposition}{Proposition}
\newtheorem{lemma}{Lemma}
\newtheorem{remark}{Remark}

\title{Shared Directional Stochastic Modes for Three-Dimensional Marine
Motion Estimation Using a Quaternion Multiplicative Extended Kalman Filter
(experimental)}
\author{%
  \IEEEauthorblockN{Mikhail S. Grushinskiy}
  \IEEEauthorblockA{Independent Researcher, USA, Fair Lawn\\
  Bareboat Necessities GitHub Project\\
  Email: mgrouch@users.github.com}%
}

\begin{document}
\maketitle

\begin{abstract}
We derive a quaternion multiplicative extended Kalman filter (QMEKF) for
the displacement, velocity and acceleration of a vessel in waves. Its
translational prior is a set of shared directional stochastic modes: on
every axis a complexified CARMA(4,1) process whose real part has exactly the
law of the independent-axis prior and whose imaginary part is its quadrature
companion. The axes share these dynamics and are driven by circular complex
noise with a Hermitian excitation, so the quadrature between horizontal and
vertical wave motion can be represented. The filter carries this law in its
minimal realization, six real states per axis and 27 error states in all,
whose coupling-off model is the independent-axis filter exactly, adaptation
included. The excitation is estimated from
the filter's own world-frame acceleration. A coupling-confidence state
machine scales it, and a covariance decoupling that provably preserves
positive semidefiniteness withdraws it. We give the continuous-time model,
its minimal realization, the exact discretization, the measurement update and error-state reset, the
coupling-confidence and fallback procedure with its positivity proof, and
stability, observability and cost arguments. We also place the model within
multivariate CARMA theory, including the complex Ornstein--Uhlenbeck
factorization of Fasen-Hartmann and Scholz and Stelzer's 2024 correction of
the MCARMA state-space characterization. The filter is experimental and is
not the shipping estimator; the companion article reports the numerical
validation.
\end{abstract}

\section{Introduction and Motivation}
The CARMA3 QMEKF of ocean-IMU models surge, sway and heave as three
independent CARMA(4,1) processes sharing one measured operating point. Linear
wave theory couples them. A wave component moves a floating hull along its
direction of travel a quarter period behind the elevation and heaves it
through the heave response amplitude operator (RAO)
\cite{longuethiggins63,faltinsen90,fossen11}. On the benchmark hull, the
$x$--$z$ cross-spectral phase at the JONSWAP peak is \SI{-62}{\degree} and
the coherence is 0.97. The coupled-prior study of this repository compared
correlated excitation (A), shared directional modes (B, ``Option C'') and a
constrained MCARMA (C) against the independent baseline. In 50-seed paired
Monte Carlo runs, only B improved both horizontal and vertical errors on
every steady directional sea with an estimated coupling. The same study
found two failures. A confidently reversed coupling gave NEES 380 against
6.8 for the correct one, and the estimated coupling lost on manoeuvring and
on wave onset. This article turns B into a standalone filter and adds the
missing safeguard.

\section{Limitations of Independent-Axis CARMA}
With independent axes, $\mat S_p(\omega)$ is diagonal: no cross-axis
information is used. Correlated excitation through identical dynamics,
$\mat Q=\mat\Sigma_w\otimes\mat Q_C$, yields
$\mat S_p(\omega)=|H(i\omega)|^2\mat\Sigma_w$, which is real at every
frequency. It can carry only the co-spectrum, 35--47\,\% of the
$x$--$z$ cross-spectral energy on single-system seas, while the physical
relation is mainly quadrature. A model that represents quadrature needs
different dynamics on the coupled axes or complex excitation. The shared
modes below are the second option, at the cost of six instead of four
states per axis (\cref{sec:minimal}).

\section{Shared Stochastic Wave-Mode Formulation}
Let $P(s)=(s^2+2\zeta_1\omega_1s+\omega_1^2)(s^2+2\zeta_2\omega_2s+\omega_2^2)$
and $B(s)=s+\omega_{z0}$ define the per-axis CARMA(4,1) transfer function
$H=gB/P$ (C41J: $f_1=\SI{0.307}{Hz}$, $\zeta_1=0.104$, $f_2=\SI{0.815}{Hz}$,
$\zeta_2=0.478$, $\omega_{z0}=0$). Write its roots as
$\mu_1,\bar\mu_1,\mu_2,\bar\mu_2$ with $\operatorname{Im}\mu_k>0$, and let
$P_+(s)=(s-\mu_1)(s-\mu_2)$.
\begin{proposition}[Complexification]
\label{prop:cx}
Let $\tfrac12(|P_+(i\omega)|^2+|P_+(-i\omega)|^2)=\omega^4+\beta\omega^2+\gamma$
and $D(s)=s^2+d_1s+d_0$ with $d_0=\sqrt\gamma$, $d_1=\sqrt{\beta+2d_0}$.
Then $D$ is Hurwitz and, for every real-coefficient $B$,
$|\tfrac{B}{P_+D}(i\omega)|^2+|\tfrac{B}{P_+D}(-i\omega)|^2=2|\tfrac{B}{P}(i\omega)|^2$.
Hence, for $\zeta=(B/P_+D)(s)\xi$ with $\xi$ circular complex white noise of
intensity $2g^2$, $\operatorname{Re}\zeta$ is a Gaussian process with
spectrum $g^2|B/P|^2$, the CARMA(4,1) process itself.
\end{proposition}
\begin{proof}
$|P(i\omega)|^2=|P_+(i\omega)|^2|P_+(-i\omega)|^2$ because the roots come in
conjugate pairs. The left side is therefore
$|B|^2(|P_+(-i\omega)|^2+|P_+(i\omega)|^2)/(|P|^2|D|^2)$.
$|D(i\omega)|^2=\omega^4+(d_1^2-2d_0)\omega^2+d_0^2$ equals the
half-sum by the choice of $d_0,d_1$, and $d_0,d_1>0$ because the half-sum is
positive for all $\omega$. For circular $\xi$, the real part has spectrum
$\tfrac14(S_\zeta(\omega)+S_\zeta(-\omega))$.
\end{proof}
For C41J the roots of $D$ are $-7.93$ and $-1.25$\,rad/s. The mirror ratio
$|H(-i\omega)|^2/|H(i\omega)|^2$, which measures how far
$\operatorname{Im}\zeta$ is from an exact Hilbert companion, is
\num{8e-4} at the peak and 0.6--2\,\% at 0.45--\SI{1}{Hz} ($\kappa=1$).

\section{Connections to Multivariate CARMA}
A $d$-variate MCARMA$(p,q)$ process is $\vct Y=\mat P(D)^{-1}\mat B(D)D\vct L$
with matrix polynomials \cite{marquardt07}. Its state-space characterization
\cite{schlemm12} was corrected in 2024 \cite{stelzer24}, which restates the
equivalence through the observer and controller canonical forms of left and
right matrix-fraction descriptions with monic denominators. We use the
corrected form. The present model has transfer matrix
$\mat H(s)=\tfrac{B(s)}{P_+(s)D(s)}\mat I_3$ applied to circular noise with
covariance $2\mat\Sigma$: a left fraction with scalar, hence commuting,
denominator, so left and right descriptions coincide.

Fasen-Hartmann and Scholz \cite{fasen21} decompose an MCARMA process with an
irreducible fraction and a complete set of regular right solvents
$\mat R_k$ into $p$ dependent complex Ornstein--Uhlenbeck processes driven
by one L\'evy process, $\vct Y=\sum_k\vct Y^{(k)}$,
$d\vct Y^{(k)}=\mat R_k\vct Y^{(k)}dt+\mat C_kd\vct L$. For one axis
($d=1$) this is the partial-fraction form $p=2\operatorname{Re}(z_1+z_2)$,
$dz_k=\mu_kz_kdt+r_kgdW$, $r_k=B(\mu_k)/P'(\mu_k)$ \cite{brockwell11}. It is
the exact modal realization of the baseline. Two of its consequences fix the
present design. First, the residues satisfy
$\sum_k\operatorname{Re}r_k=\sum_k\operatorname{Re}(r_k\mu_k)=0$ (the first
two Markov parameters of $B/P$ vanish), so independently driven oscillators
lose relative degree three. They give white velocity and an acceleration
with no finite variance. Second, per-mode complex gains $\beta_k$ on a
second axis keep relative degree three iff
$\sum\operatorname{Re}(\beta_kr_k)=\sum\operatorname{Re}(\beta_kr_k\mu_k)=0$,
that is, iff that axis is $(b_1B+b_0)/P$ with real $b_0,b_1$. Its phase
$-\arctan(b_0/b_1\omega)$ is never constant. The four-state modal form
therefore cannot carry a constant cross-axis quadrature. The complexified mode
of \cref{prop:cx} adds the two real poles of $D$ and can, at relative degree
three by construction.

An unrestricted MCARMA(4,1) on three axes has 54 coefficients. From
\SI{20}{min} of three-axis acceleration their Whittle Fisher information has
rank 51, and 44\,\% of them have a relative Cram\'er--Rao bound above
100\,\% \cite{whittle53,fasen22}. The shared-mode model has 10
identifiable parameters (frequency scale, three amplitudes, six cross
terms). It is preferred here because its structure is the physics, a
single directional wave field seen on three axes, and because it is
identifiable online and costs a fixed 18 motion states.

\section{Continuous-Time Stochastic State-Space Model}
Per axis $a\in\{x,y,z\}$ the complex kinematic state is
$\vct\zeta_a=[\zeta,\dot\zeta,\ddot\zeta,l]\T$, $\dot{\vct\zeta}_a=\mat A_c\vct\zeta_a+\vct e_3\xi_a$,
with $\mat A_c$ the complex companion form of $P_+D$ in these coordinates,
$\zeta=\vct c\T\vct\zeta$ and the noise entering $\ddot\zeta$ only. The real
state is $\vct x_a=[\operatorname{Re}\vct\zeta_a;\operatorname{Im}\vct\zeta_a]\in\R^8$,
$\mat A_8=\left[\begin{smallmatrix}\operatorname{Re}\mat A_c&-\operatorname{Im}\mat A_c\\
\operatorname{Im}\mat A_c&\operatorname{Re}\mat A_c\end{smallmatrix}\right]$,
and the three axes share it:
\begin{equation}
d\vct x_L=(\mat I_3\otimes\mat A_8)\vct x_Ldt+d\vct\beta,\quad
E[d\vct\beta d\vct\beta\T]=\mat W\,dt,
\label{eq:sde}
\end{equation}
\begin{equation}
\mat W_{ab}=g^2\bigl(\operatorname{Re}(\Sigma_{ab})\mat E+\operatorname{Im}(\Sigma_{ab})\mat J\mat E\bigr),
\label{eq:W}
\end{equation}
with $\mat E=\operatorname{diag}(\vct e_3\vct e_3\T,\vct e_3\vct e_3\T)$,
$\mat J=\left[\begin{smallmatrix}\mat 0&-\mat I_4\\\mat I_4&\mat 0\end{smallmatrix}\right]$.
\begin{lemma}
\label{lem:W}
$\mat W\succeq0$ iff $\mat\Sigma\succeq0$, and $\mat J$ commutes with $\mat A_8$
and $\mat E$.
\end{lemma}
\begin{proof}
$\mat J$ is multiplication by $i$ in the real representation, so it commutes
with the realification of any complex matrix ($\mat A_8$, $\mat E$). For
$\vct u\in\R^{24}$ write $u_a=\vct u_a^{\mathrm{Re}}+i\vct u_a^{\mathrm{Im}}$.
Then $\vct u\T\mat W\vct u=g^2\sum_{a,b}\operatorname{Re}(\bar w_a\Sigma_{ab}w_b)$
with $w_a$ the third component of $u_a$, which is
$g^2\operatorname{Re}(\vct w\Hr\mat\Sigma\vct w)\ge0$. Conversely, $\vct w$
ranges over $\C^3$.
\end{proof}
The displacement readouts are $p_a=\vct e_1\T\vct x_a$ (real part of
$\zeta$), $v_a=\vct e_2\T\vct x_a$, $a_a=\vct e_3\T\vct x_a$. These are unit
selectors, and $\dot v_a=a_a$ holds without noise (relative degree three
preserved). The displacement cross-spectral matrix is
\begin{equation}
S_{ab}(\omega)=\tfrac{g^2}{2}\bigl[\Sigma_{ab}|H(i\omega)|^2+\bar\Sigma_{ab}|H(-i\omega)|^2\bigr],
\label{eq:S}
\end{equation}
with $H=B/(P_+D)$. Away from the small mirror term, $\arg S_{ab}=\arg\Sigma_{ab}$
and the coherence is $|\Sigma_{ab}|^2/(\Sigma_{aa}\Sigma_{bb})$. A real
$\mat\Sigma$ reproduces correlated excitation in law; a real diagonal one the
independent baseline.

\section{Minimal Realization}
\label{sec:minimal}
The realified mode carries two directions per axis that no measurement sees.
Each root $r$ of $D$ is an eigenvalue of both $\mat A_c$ and $\bar{\mat A}_c$,
so $\mat A_8$ has it twice with a two-dimensional eigenspace, which the
scalar output $\operatorname{Re}\ddot\zeta$ cannot resolve. The unobservable
subspace is
\begin{equation}
\mathcal U=\{\vct 0\}\times\bigl(\operatorname{span}_\C\{\vct w(r_3),\vct w(r_4)\}\cap\R^4\bigr),\quad
\vct w(r)=[B(r),rB(r),r^2B(r),1]\T,
\end{equation}
two real directions in the imaginary half, independent of $\mat\Sigma$, of the
attitude and of every QMEKF operation. The observable quotient has six
states per axis and the same law of $p,v,a$ on all three axes, every co- and
quadrature spectrum included. The filter uses the following realization of
it. With $P_+=N_r+iN_i$ ($N_r,N_i$ real polynomials),
\begin{equation}
\operatorname{Re}\zeta=\frac BP\,\eta,\quad\eta=\mat V\begin{bmatrix}\xi_r\\\xi_i\end{bmatrix},\quad
\mat V=\frac{[N_r,\;N_i]}D,
\label{eq:cascade}
\end{equation}
and $|D|^2=|N_r|^2+|N_i|^2$ makes $\mat V$ lossless (co-inner). The six states
of an axis are the baseline's own $[p,v,a,l]$, driven at $a$ by
$\eta=\xi_r+\mat C_w\vct w$, and the two states $\vct w$ of $\mat V$ in
input-normal coordinates: $\mat A_w+\mat A_w\T=-\mat B_w\mat B_w\T$,
$\mat C_w=-\mat B_w\vct e_1\T$, $\mat B_w$ upper triangular. Its noise
input is $\mat G=[\vct e_3,\vct 0;\mat B_w]$, and the Hermitian excitation
enters through two intensities,
\begin{equation}
\mat W_{ab}=\operatorname{Re}(\Sigma_{ab})\mat W_R+\operatorname{Im}(\Sigma_{ab})\mat W_I,\quad
\mat W_R=g^2\mat G\mat G\T,\ \mat W_I=g^2\mat G\mat J_2\mat G\T,
\label{eq:W6}
\end{equation}
$\mat J_2=\left[\begin{smallmatrix}0&-1\\1&0\end{smallmatrix}\right]$, $\mat W_I$
antisymmetric. Every quantity scales exactly with $\kappa$
($\mat V_\kappa(s)=\mat V_1(s/\kappa)$), so one closed-form $2\times2$
construction at $\kappa=1$ is the only factorization. The derivations, the
minimal stochastic order of the three-axis law ($4\operatorname{rank}\mat\Sigma
+2\operatorname{rank}\operatorname{Im}\mat\Sigma$) and the reasons for this
realization are in the companion article on the minimal realization.

\section{State-Vector Definition}
The error state of the QMEKF is
\begin{equation}
\delta\vct x=[\delta\vct\theta,\;\delta\vct b_g,\;\delta\vct x_x,\;\delta\vct x_y,\;\delta\vct x_z,\;\delta\vct b_a]\in\R^{27},
\end{equation}
with $\delta\vct\theta$ the body-frame attitude error of
$\vct q=\hat{\vct q}\otimes\exp(\delta\vct\theta/2)$, the gyroscope and
accelerometer bias errors, and the 18 motion states
$\vct x_a=[p_a,v_a,a_a,l_a,\vct w_a\T]\T$ of \cref{sec:minimal} in NED with $z$ down. The
motion states are additive; they are world-frame quantities, because a wave
field is fixed in the world.

\section{QMEKF Attitude and Bias States}
These are the baseline's. Gyroscope bias is a random walk,
$\dot{\vct b}_g=\vct w_{bg}$; the accelerometer bias is a random walk with
temperature compensation. The attitude kinematics is
$\dot{\vct q}=\tfrac12\vct q\otimes(\vct\omega_m-\vct b_g-\vct n_g)$, and its
linearized error dynamics are
$\delta\dot{\vct\theta}=-[\hat{\vct\omega}]_\times\delta\vct\theta-\delta\vct b_g-\vct n_g$,
with $\hat{\vct\omega}=\vct\omega_m-\hat{\vct b}_g$ \cite{markley03}.

\section{Coupling and Cross-Spectral Relationships}
For a linear wave component of azimuth $\psi$ and elevation $\eta$, the hull
moves $[-i\,s_N,-i\,s_E,H_z(\omega)]\eta$ (surge and sway along the
direction of travel in quadrature, heave through its RAO). The regression of
a horizontal axis on heave at the operating frequency is therefore complex:
$\beta_i=S_{iz}/S_{zz}$. With the axes' accelerations
$\vct a=-\omega^2\vct p$, the demodulated band cross-spectrum of the
accelerations has the same ratios. The excitation that reproduces the
measured regression at $f_{op}$ is
\begin{equation}
\mat\Sigma=\begin{bmatrix}\sigma_h^2\mat I_2&\vct b\\\vct b\Hr&\sigma_z^2\end{bmatrix},\quad
\vct b=s\,c\,\kappa_c\frac{\sigma_z^2}{\hat S_{zz}}\hat{\vct S}_{hz},
\label{eq:map}
\end{equation}
with $s=\max(0,1-1/(N_{\mathrm{eff}}\gamma^2))$ the debiased-coherence
shrinkage, $\kappa_c$ the configured coupling gain, $c$ the coupling
confidence (\cref{sec:conf}). If
$q=\vct b\Hr(\sigma_h^2\mat I_2)^{-1}\vct b>\rho\sigma_z^2$ ($\rho=0.8$), the
horizontal block is enlarged by $q/(\rho\sigma_z^2)$, keeping the measured
regression (an earlier policy scaled $\vct b$ down to $0.999\sigma_z^2$). The
Schur complement of the horizontal block is then nonnegative, so
$\mat\Sigma\succeq0$ by construction. Every factor is invariant under rotations about the
vertical, so the prior does not depend on how the waves sit relative to
north. The direction of travel follows from $\vct\beta=e^{i\phi}r\vct u$:
$\phi=\tfrac12\arg(\beta_x^2+\beta_y^2)$, $\vct u=\operatorname{Re}(e^{-i\phi}\vct\beta)$,
with the sign of $\vct u$ chosen so that the along-direction phase lies in
$(-\pi,0)$ (the quadrature lag). Its large-sample uncertainty is
$\sqrt{(1-\gamma^2)/(2N\gamma^2)}$.

\section{Process-Noise Covariance}
The full process noise is block-diagonal across the attitude, bias and
motion groups. The motion block is $\mat W$ of \eqref{eq:W6}. Sensor noise
(accelerometer, gyroscope, magnetometer white noise) enters only the
measurement and the attitude propagation. The wave excitation $\mat\Sigma$
models the hull's motion, not the sensor, so it is not double-counted: the
envelope estimate uses the filter's own world-frame acceleration, whose
sensor-noise content in the band is $\sigma_n^2\,dt\,\pi bf/2$ ($b$ the band
fraction). The coherence shrinkage removes it from the cross terms, and the
diagonal keeps the wrapper's amplitude law, which has its own noise
handling.

\section{Exact Discrete-Time Propagation}
With $\mat A_6$ the dynamics of one axis, $\mat F_6=e^{\mat A_6h}$ and
$\mat Q_X=\int_0^he^{\mat A_6t}\mat W_Xe^{\mat A_6\T t}dt$ for $X\in\{R,I\}$, the
exact discretization is
\begin{equation}
\mat F_L=\mat I_3\otimes\mat F_6,\quad
\mat Q_{L,ab}=\operatorname{Re}(\Sigma_{ab})\mat Q_R+\operatorname{Im}(\Sigma_{ab})\mat Q_I,
\label{eq:disc}
\end{equation}
$\mat Q_R$ symmetric and $\mat Q_I$ antisymmetric. Both are computed by one
paired scaled Taylor series and doublings (the baseline's, exact up to a
controlled truncation below float rounding), and by its Taylor extension on a
jittered clock. $\mat Q_L\succeq0$ follows from \cref{lem:W}: the six-state
block is the observable quotient of \eqref{eq:sde}, and its noise is the
image of $\mat W$. The dynamics are stable: the poles of $\mat A_6$ are
$\mu_1,\bar\mu_1,\mu_2,\bar\mu_2$ and the roots of $D$, all in the open left
half-plane for every $\kappa>0$.

\section{Measurement Model}
The accelerometer measures the specific force in the body frame,
\begin{equation}
\vct f=\mat R_{wb}(\vct a_w-g\vct e_z)+\vct b_a+\vct n_a,
\end{equation}
$\vct a_w=[a_x,a_y,a_z]\T$ the motion readouts. Its Jacobian has
$\partial\vct f/\partial\delta\vct\theta=[\mat R_{wb}(\hat{\vct a}_w-g\vct e_z)]_\times$,
$\partial\vct f/\partial\delta\vct b_a=\mat I$, and on the motion block the
columns $\mat R_{wb}\vct e_i$ at the $a$ selector of axis $i$. The
magnetometer measures $\mat R_{wb}\vct m_w$ and corrects attitude only,
through the wrapper's service gates.

\section{Kalman Update Equations}
With $\mat S=\mat H\mat P\mat H\T+\mat R$, $\mat K=\mat P\mat H\T\mat S^{-1}$
and the innovation $\vct r$, the error mean is $\delta\hat{\vct x}=\mat K\vct r$.
The covariance follows the Joseph form
$\mat P^+=(\mat I-\mat K\mat H)\mat P(\mat I-\mat K\mat H)\T+\mat K\mat R\mat K\T$,
which is positive semidefinite for any gain. The motion and bias states are
corrected additively. The NIS $\vct r\T\mat S^{-1}\vct r$ is reported.

\section{Error-State Reset}
After each update the attitude error is folded into the quaternion,
$\hat{\vct q}\leftarrow\hat{\vct q}\otimes\exp(\delta\hat{\vct\theta}/2)$,
and the covariance is transported with the reset Jacobian
$\mat G=\operatorname{blockdiag}(\mat I-\tfrac12[\delta\hat{\vct\theta}]_\times,\mat I)$,
$\mat P\leftarrow\mat G\mat P\mat G\T$. Only the attitude rows and columns
change. The motion states are not affected by the reset, so the
cross-covariances between attitude and motion transform with the attitude
rows exactly as in the baseline.

\section{Sea-State Adaptation}
Every \SI{0.1}{s} the wrapper commits the frequency scale
$\kappa=f_{\mathrm{tune}}/f_{z,\mathrm{base}}$ and isotropic amplitudes from
OU-III's tuner. The motion block then rebuilds $\mat A_6(\kappa)$,
$\mat F_6,\mat Q_R,\mat Q_I$ and $\mat\Sigma$. A commit holds $[p,v,a,l]$,
as the baseline does, and the normalized colouring states
$\vct w_a/s_a$, $s_a=\sqrt{g^2\sigma_a^2}$: $\vct w_a\leftarrow(s_a'/s_a)\vct w_a$ and
the covariance transforms with the same diagonal congruence. $\mat V$ is
lossless with a constant Gramian, so $\eta$ stays exactly white under any
change of $\kappa$ and $\sigma$: with the coupling off the motion block
\emph{is} the independent-axis CARMA(4,1) filter in law, adaptation
included. The covariance is never reset during adaptation (only on entering
Live, to the stationary prior), so a commit changes the next prediction, not
the posterior.

\section{Coupling-Confidence Estimation}
\label{sec:conf}
Let $\vct\beta_p$ be the regression of the prior (the \SI{240}{s} estimate, or
an external prior) and $\vct\beta_f$ that of a \SI{60}{s} monitor fed with the
same world-frame accelerations. For an estimate from $N$ independent
envelope samples, $\operatorname{var}\hat\beta_i\approx(1-\gamma_{iz}^2)S_{ii}/(NS_{zz})$
\cite{bendat10}. The mismatch statistic is
\begin{equation}
D^2=\sum_{i\in\{x,y\}}\frac{|\beta_{f,i}-\beta_{p,i}|^2}{v_{f,i}+v_{p,i}+\epsilon^2|\beta_{p,i}|^2}.
\end{equation}
For a consistent prior and independent estimates $D^2$ is approximately
$\chi^2_4/2$-distributed, with $E[D^2]=2$. The overlap of the two estimates
only lowers it. Hysteresis declares the prior inconsistent for $D^2>16$ and
consistent again for $D^2<6$. Further evidence gates are the monitor's SNR
against the envelope noise floor (log-linear between 10 and 100), the
low-passed world yaw rate (manoeuvre, $>\SI{3}{\degree/s}$), the low-passed
$|d\ln f_{op}/dt|$ (sea change, $>0.01/$s), both held \SI{30}{s}, and the
ratio of the \SI{20}{s} to the \SI{600}{s} mean NIS ($>2$ fallback, $<1.5$
recovery). Their product $e\in[0,1]$ drives
\begin{equation}
c\leftarrow c+(1-e^{-dt/\tau})(e-c),\quad\tau=\begin{cases}\SI{3}{s},&e<c,\\\SI{60}{s},&e\ge c,\end{cases}
\end{equation}
so $|\dot c|\le1/\tau$: withdrawal is fast, restoration slow, and there is no
rapid switching.

\section{Safe Decoupling and Fallback}
\label{sec:fallback}
The confidence enters \eqref{eq:map}, so $c=0$ gives the real diagonal
excitation, which is the independent-axis CARMA(4,1) law on every axis. The
state is not modified, so position, velocity and attitude are continuous.
The covariance, however, still carries horizontal--vertical correlation built
by the coupled prior. Partition the error state into $\vct r$ (attitude and
biases, 9 states) and the motion states $\vct m=(\vct h,\vct z)$ (12 + 6).
\begin{proposition}[Positive decoupling]
\label{prop:shrink}
Let $\mat P\succeq0$ with $\mat P_{rr}\succ0$,
$\mat E=\mat P_{hr}\mat P_{rr}^{-1}\mat P_{rz}$ and $\alpha\in[0,1]$. Replacing
$\mat P_{hz}$ by $\mat E+\alpha(\mat P_{hz}-\mat E)$ (and $\mat P_{zh}$ by its
transpose) leaves $\mat P$ positive semidefinite.
\end{proposition}
\begin{proof}
$\mat P\succeq0$ iff $\mat P_{rr}\succ0$ and the Schur complement
$\mat C=\mat P_{mm}-\mat P_{mr}\mat P_{rr}^{-1}\mat P_{rm}\succeq0$. The
replacement leaves $\mat P_{rr},\mat P_{rm}$ and the diagonal blocks of
$\mat P_{mm}$ unchanged and changes $\mat C_{hz}=\mat P_{hz}-\mat E$ into
$\alpha\mat C_{hz}$. Hence $\mat C'=\alpha\mat C+(1-\alpha)\operatorname{blockdiag}(\mat C_{hh},\mat C_{zz})$.
$\operatorname{blockdiag}(\mat C_{hh},\mat C_{zz})\succeq0$ because principal
submatrices of $\mat C\succeq0$ are, and a convex combination of PSD matrices
is PSD.
\end{proof}
The filter applies this with $\alpha=c_{\mathrm{new}}/c_{\mathrm{old}}$
whenever $c$ falls (at most at \SI{10}{Hz}). Simply zeroing $\mat P_{hz}$
would not be safe. In general the matrix with $\mat P_{hz}=\mat 0$ and
nonzero $\mat P_{hr},\mat P_{rz}$ is indefinite. At $\alpha=0$ the
correlation explained by the shared attitude and bias errors, $\mat E$, is
kept and the rest removed. The dynamics are block-diagonal across axes, so
with $\mat Q_{hz}=\mat 0$ any remaining cross-covariance decays at the modes'
damping rates. Recovery is the reverse: when the evidence returns, $c$
rises with \SI{60}{s} and the cross-covariance rebuilds through
\eqref{eq:disc} and the updates.

\section{Stability and Observability}
For fixed parameters the prediction is the exact transition of a stable
linear system \eqref{eq:disc}. The accelerometer observes
the three $a$ readouts and, through gravity, two tilt angles; the
magnetometer observes heading. Within the motion block, each axis is a
minimal realization of the transfer function \eqref{eq:cascade}
(McMillan degree six), hence observable from its acceleration readout for
every $\mat\Sigma$; the directions $\mathcal U$ that the realified mode could
not resolve are not carried at all. The pair is therefore detectable for
every $\mat\Sigma$, and the Riccati recursion has a bounded solution for fixed
parameters \cite{andersonmoore79}. Two
qualifications apply. Displacement is observable only through the prior's
low-frequency behaviour (the zero of $B$ at DC), as for the baseline. With
complex $\mat\Sigma$ the $x$--$z$ quadrature also makes heading weakly
observable from the accelerometer, which is information when the coupling is
right and a bias when it is wrong. That is the reason for \cref{sec:conf}.
This is a linear-model argument at frozen parameters and frozen attitude
linearization; it is not a proof of stability of the adaptive nonlinear
filter.

\section{Computational Complexity}
The prediction costs the attitude and bias blocks plus the motion block:
$6\times6$ congruences per axis pair (block-diagonal $\mat F_L$), the
structured $\mat Q_L$ \eqref{eq:disc}, and the $6\times18$ attitude--motion
cross-covariances. With the accelerometer update, a Joseph update with a
three-row measurement on 27 states, a sample costs about 8\,000
multiply-adds against 4\,200 for the baseline. The confidence adds two envelope estimators
(36 multiply-adds per sample) and, at most at \SI{10}{Hz} while $c$ falls, one
$9\times9$ LDLT and a $12\times9\times6$ product. The model rebuild
(\SI{10}{Hz}) computes $\mat F_6,\mat Q_R,\mat Q_I$ once for all axes. All
matrices are fixed size, nothing is allocated per update, and the 27-state
QMEKF is one heap allocation at start-up. Measured sizes and times are in the
companion article.

\section{Limitations and Future Research}
The coupling is estimated in one band at one operating frequency. Its phase
is constant across frequency, while the true $x$--$z$ phase follows the heave
RAO. Two-system seas with different directions are represented only by the
rank of $\mat\Sigma$, and a two-mode extension needs a second-peak frequency
estimator that does not exist yet. The confidence detects only mismatches
that the IMU's own band statistics can contradict, and it reacts after
seconds to a minute of evidence. The complex mode costs two extra states per
axis. A constrained MCARMA with the heave RAO in the transfer matrix gave
larger gains in the study but needs hull knowledge and more stack, and is not
pursued here. Hardware runtime and stack high-water marks remain to be
measured on the AtomS3R.

\begin{thebibliography}{99}
\bibitem{longuethiggins63}
M. S. Longuet-Higgins, D. E. Cartwright and N. D. Smith, ``Observations of
the directional spectrum of sea waves using the motions of a floating buoy,''
in \emph{Ocean Wave Spectra}, Prentice-Hall, 1963, pp. 111--136.
\bibitem{faltinsen90}
O. M. Faltinsen, \emph{Sea Loads on Ships and Offshore Structures}.
Cambridge University Press, 1990.
\bibitem{fossen11}
T. I. Fossen, \emph{Handbook of Marine Craft Hydrodynamics and Motion
Control}. Wiley, 2011.
\bibitem{marquardt07}
T. Marquardt and R. Stelzer, ``Multivariate CARMA processes,''
\emph{Stochastic Process. Appl.}, vol. 117, no. 1, pp. 96--120, 2007.
\bibitem{schlemm12}
E. Schlemm and R. Stelzer, ``Multivariate CARMA processes, continuous-time
state space models and complete regularity of the innovations of the sampled
processes,'' \emph{Bernoulli}, vol. 18, no. 1, pp. 46--63, 2012.
\bibitem{stelzer24}
R. Stelzer, ``Correction to: Multivariate CARMA processes, continuous-time
state space models and complete regularity of the innovations of the sampled
processes,'' arXiv:2411.06935, 2024.
\bibitem{fasen21}
V. Fasen-Hartmann and M. Scholz, ``Factorization and discrete-time
representation of multivariate CARMA processes,'' arXiv:2102.11681, 2021.
\bibitem{brockwell11}
P. J. Brockwell, R. A. Davis and Y. Yang, ``Estimation for non-negative
L\'evy-driven CARMA processes,'' \emph{J. Bus. Econom. Statist.}, vol. 29,
pp. 250--259, 2011.
\bibitem{whittle53}
P. Whittle, ``The analysis of multiple stationary time series,'' \emph{J.
Roy. Statist. Soc. B}, vol. 15, pp. 125--139, 1953.
\bibitem{fasen22}
V. Fasen-Hartmann and C. Mayer, ``Whittle estimation for continuous-time
stationary state space models with finite second moments,'' \emph{Ann. Inst.
Statist. Math.}, vol. 74, pp. 233--270, 2022.
\bibitem{markley03}
F. L. Markley, ``Attitude error representations for Kalman filtering,''
\emph{J. Guid. Control Dyn.}, vol. 26, no. 2, pp. 311--317, 2003.
\bibitem{bendat10}
J. S. Bendat and A. G. Piersol, \emph{Random Data: Analysis and Measurement
Procedures}, 4th ed. Wiley, 2010.
\bibitem{andersonmoore79}
B. D. O. Anderson and J. B. Moore, \emph{Optimal Filtering}. Prentice-Hall,
1979.
\end{thebibliography}
\end{document}
